\cotefile{69}
\scanpage{68}
\diagram{\(\begin{array}{@{}r@{\ }l@{\quad}l}
 &\mathrm{Pl}(\unsure{L|D}_1,C)=C&27=3^3\\
 \uparrow&\mathrm{Pl}(L_2,C)&27.16=2^4 3^3=432\\
 \uparrow&\mathrm{Pl}(L_3,C)&27.16.10=2^5.3^3.5=4320\\
 \uparrow&\mathrm{Pl}(L_4,C)\xleftarrow{\sim}\mathrm{Pl}''(L_5,C)&27.16.10.6=2^6.3^4.5=25.920\\
 \unsure{\uparrow|\vert}&\mathrm{Pl}'(L_5,C)\unsure{:=|:}&27.16.10.6.2=2^7.3^4.5=51840\\
 2{\uparrow}&\mathrm{Pl}'(L_6,C&,5
\end{array}\)}

On trouve pour les configurations correspondantes

\diagram{\(\begin{array}{@{}l@{\quad}l@{\quad}l}
\deleted{\unsure{1|a}} B_1(C)&27:1=27=3^3&\mathrm{aut}\,L_1=e\\
B_2(C)&27.16:2!=27.8=2^3 3^3=216&\mathrm{aut}\,L_2\unsure{:|=}\unsure{\mathfrak S|\mathscr S}_2\\
B_3(C)&27.16.10:3!=2^{\deleted{\unsure{5|3}}4}.3^2.5=720=6!&\mathrm{aut}\,L_3=\mathfrak S_3\\
B_4(C)&27.16.10.6:4!=2^3.3^3.5=1080&\mathrm{aut}\,L_4=\mathfrak S_4\\
\quad\unsure{\simeq|\approx} B_4''(C)&&\\
B'_5(C)&27.16.10.6.2:5!=2^4.3^3=432&\mathrm{aut}\,L_5=\mathfrak S_5\\
B(C)=B_6(C)&27.16.10.6.2.1:6!=2^3 3^2=72&\mathrm{aut}\,L_5\simeq\mathfrak S_6\\
\mathrm{Bib}(C)&36&\mathrm{aut}(L'_5)\simeq\mathfrak S_2\times\mathfrak S_6
\end{array}\)}

\begin{center}\rule{1.5cm}{.5pt}\\[-.8ex]\rule{1.1cm}{.4pt}\end{center}

\uline{Étude cas par cas}

1) $\mathrm{Pl}(L_1,C)=C$ \inserted{de cardinal $27$}. Si $x\in C$, $\mathrm{Aut}(C;x)$ est transitif sur l'ens \deleted{\unsure{des|de}} $C-\Xi(x)$ \inserted{de cardinal $16$} \deleted{[\unsure{des pts} $z\ne x$ \unsure{liés|à $x$}]}.

2) $\mathrm{Pl}(L_2,C)$

\uline{Prop.} Pour \unsure{deux|des|tout} couples $(x,y)$ de sommets dist. non liés, $\exists!$ base $b$ telle que $x\in b$, $y\in b'$ \inserted{base opposée} (et $x,y$ \deleted{liés} correspondent $y=\sigma_b\unsure{x=\sigma_{b'}x|x\Rightarrow\sigma_{b'}x}$, $x=\sigma_b y=\sigma_{b'}y$)

Ainsi l'ens des \deleted{couples} \inserted{parties gauches 2-libres} $(x,y)$ \deleted{de sommets} est en corr. $1$--$1$ avec l'ens des couples d'une base $b$ et d'un $x\in b$. La cat. des $(C,(x,y))$ équivaut à celle des ens. de cardinal $5$ (\unsure{ou encore}, des ens de cardinal $6$ à pt marqué)

On peut écrire $(x,y)=(\unsure{\delta_6|\delta_\sigma|J_\sigma|\xi_6},\delta'_6)$. L'ens $C-\Xi(\delta_6)-\unsure{\Sigma|\Xi|E}(\delta'_6)$ \inserted{$-\{x\unsure{+|,}y\}$} des pts \inserted{$\ne x,y$} non liés, ni à $y$ est l'ens des $\delta_{ij}$ $(1\leq i<j\leq5)$, de cardinal $10$ \inserted{$=\binom{5}{2}$}, et $\mathrm{Aut}(C,(x,y))=\mathfrak S_5$ y opère transitivement.
% archive sheet number: 67
